\subsection{Point measures. Measures with finite support}

PROPOSITION 12. --- Let \(a_i\) ( \(1 \leqslant i \leqslant n\) ) be distinct points in a locally compact space X. Every measure on X whose support is contained in the set of the \(a_i\) is a linear combination of the measures \(\varepsilon_{a_i}\) ( \(1 \leqslant i \leqslant n\) ).

For, such a measure \(\mu\) is zero for every function \(f\in\mathcal{K}(\mathrm{X};\mathbf{C})\) satisfying the \(n\) relations \(f(a_{i})=0\) (No.~3, Prop. 8); since these relations may be written \(\varepsilon_{a_{i}}(f)=0\) , \(\mu\) is a linear combination of the \(\varepsilon_{a_{i}}\) (A, II, §7, No.~5, Cor. 1 of Th. 7).

In particular, every measure whose support is either empty or reduced to a single point x is of the form \(\alpha\varepsilon_{x}\) , where \(\alpha\) is a complex number; such a measure is said to be a point measure; thus every measure whose support is finite is a sum of point measures.

THEOREM 1. --- Every measure \(\mu\) on a locally compact space \(X\) is in the vague closure of the vector space \(V\) of measures whose support is finite and contained in \(\operatorname{Supp}(\mu)\) .

It suffices to prove that \(\mu\) is orthogonal to the subspace \(V^{\circ}\) of \(\mathcal{K}(X; \mathbf{C})\) orthogonal to \(V\) (TVS, II, §6, No.~3, Cor. 2 of Th. 1), that is, that the relations \(\langle f, \varepsilon_{a} \rangle = 0\) , where \(a\) runs over the support of \(\mu\) , imply \(\langle f, \mu \rangle = 0\) ; but this is just Prop. 8 of No.~3.

COROLLARY 1. --- Every bounded measure \(\mu\) on X is in the vague closure of the convex set A of measures whose support is finite and is contained in that of \(\mu\) , and whose norm is \(\leqslant \| \mu \|\) . Moreover, if \(\nu\) tends vaguely to \(\mu\) while remaining in A, then \(\| \nu \|\) tends to \(\| \mu \|\) .

To prove the first assertion, it suffices to establish that the measure \(\mu\) belongs to the polar set of the polar set \(A^{\circ}\) of A in \(\mathcal{K}(X;C)\) (TVS, II, §6, No.~3, Th. 1 and §8, No.~4); this means that for \(f \in \mathcal{K}(X;C)\) , the relations \(|\langle f,\varepsilon_{a}\rangle| \leqslant 1/\|\mu\|\) for all \(a \in \operatorname{Supp}(\mu)\) imply that \(|\langle f,\mu\rangle| \leqslant 1\) ; but this is a consequence of Cor. 3 of Prop. 8 of No.~3.

To prove the second assertion, we note that

\[
\liminf _ {\nu \to \mu ,   \nu \in A} \| \nu \| \geqslant \| \mu \|
\]

since the function \(\nu \mapsto \| \nu \|\) is lower semi-continuous for the vague topology (§1, No.~9, Cor. 4 of Prop. 15), and the conclusion follows from the fact that \(\| \nu \| \leqslant \| \mu \|\) for \(\nu \in A\) by definition.

COROLLARY 2. --- Every bounded measure \(\mu\) on \(X\) is in the vague closure of the set of measures whose support is finite and contained in that of \(\mu\) and whose norm is equal to \(\|\mu\|\) .

We can suppose that \(\mu \neq 0\) . Let \(V\) be an open neighborhood of 0 for the vague topology; for every \(\varepsilon\) such that \(0 < \varepsilon < 1\) , there exists, by virtue of Cor. 1, a measure \(\nu_0\) whose support is finite and contained in \(\operatorname{Supp}(\mu)\) and for which \(\nu_0 - \mu \in V\) and \(\| \mu \| \geqslant \| \nu_0 \| \geqslant (1 - \varepsilon) \| \mu \|\) . Setting \(\nu = (\| \mu \| / \| \nu_0 \|) \nu_0\) , we have \(\| \nu \| = \| \mu \|\) and \(\| \nu - \nu_0 \| \leqslant \| \mu \|\) ; for \(\varepsilon\) sufficiently small we therefore have \(\nu - \mu \in V + V\) , whence the conclusion.

COROLLARY 3. --- Every bounded positive measure \(\mu\) on \(X\) is in the vague closure of the convex set of positive measures whose support is finite and contained in that of \(\mu\) and whose norm is equal to \(\|\mu\|\) .

The same reasoning as in Cor. 2 shows that we can limit ourselves to proving that \(\mu\) is in the vague closure of the convex set B formed by the positive measures with finite support contained in \(\mathrm{Supp}(\mu)\) and with norm \(\leqslant \| \mu \|\) . Again, it suffices to establish that \(\mu\) belongs to the polar set of \(\mathbf{B}^{\circ}\) , the polar set of B in the space \(\mathcal{K}(\mathbf{X};\mathbf{R})\) (TVS, II, §6, No.~3, Th. 1); but this means that for \(f\in \mathcal{K}(\mathbf{X};\mathbf{R})\) the relations \(\langle f,\varepsilon_a\rangle \geqslant -1 / \| \mu \|\) for all \(a\in \mathrm{Supp}(\mu)\) imply \(\langle f,\mu \rangle \geqslant -1\) , which is a consequence of No.~3, Cor. 2 of Prop. 8.

COROLLARY 4. --- In the space \(\mathcal{M}(\mathrm{X};\mathbf{C})\) , the set of point measures is total for the topology of strictly compact convergence (§1, No.~10).

On the cone \(\mathcal{M}_{+}(X)\) , the topology of strictly compact convergence is identical to the vague topology (§1, No.~10, Prop. 18), and every measure on X may be written \(\mu_{1}-\mu_{2}+i\mu_{3}-i\mu_{4}\) , where the \(\mu_{j}\ (1\leqslant j\leqslant4)\) are positive measures; the conclusion therefore follows from Th. 1.

PROPOSITION 13. --- Let \(\mu\) be a measure on a locally compact space \(X\) . For a point \(x_0\) to belong to \(\operatorname{Supp}(\mu)\) , it is necessary and sufficient that the point measure \(\varepsilon_{x_0}\) be in the vague closure of the set of measures \(g \cdot \mu\) , where \(g\) runs over the set of continuous functions with compact support such that \(\| g \cdot \mu \| \leqslant 1\) .

The condition is obviously sufficient by virtue of Prop. 6 of No.~2. To see that it is necessary, suppose \(x_0 \in \operatorname{Supp}(\mu)\) ; consider a finite number of functions \(f_k\) ( \(1 \leqslant k \leqslant n\) ) in \(\mathcal{K}(X; \mathbf{C})\) , and an arbitrary number \(\delta > 0\) ; we are to prove that there exists a function \(g \in \mathcal{K}(X; \mathbf{C})\) such that \(\| g \cdot \mu \| \leqslant 1\) and such that

\[
\left| f _ {k} \left(x _ {0}\right) - \mu \left(g f _ {k}\right) \right| \leqslant \delta
\]

for \(1 \leqslant k \leqslant n\) . Let U be a relatively compact open neighborhood of \(x_0\) such that the oscillation of each of the \(f_k\) ( \(1 \leqslant k \leqslant n\) ) on U is \(\leqslant \delta/2\) . By hypothesis, since \(x_0 \in \text{Supp}(\mu)\) , there exists a function \(g_0 \in \mathcal{K}(X; \mathbf{C})\) whose support is contained in U and such that \(\mu(g_0) \neq 0\) ; the measure \(\nu = g_0 \cdot \mu\) is not zero, since for every function \(f \in \mathcal{K}(X; \mathbf{C})\) equal to 1 on U, \(\nu(f) = \mu(g_0) \neq 0\) . Moreover, \(\nu\) is bounded (No.~3, Prop. 11);

multiplying \(g_{0}\) by a scalar, we can suppose that \(\|\nu\|=1\) . This being so, setting \(\alpha_{k}=f_{k}(x_{0})\) we can write, for \(1\leqslant k\leqslant n\) and for every function \(h\in\mathcal{K}(\mathrm{X};\mathbf{C})\) ,

\[
f _ {k} (x _ {0}) - \nu (f _ {k} h) = \alpha_ {k} (1 - \nu (h)) + \nu ((\alpha_ {k} - f _ {k}) h).
\]

Since \(\nu\) has its support in U, we may identify it with its restriction to U; the hypothesis \(\| \nu \| = 1\) then implies that there exists a function \(h \in \mathcal{K}(\mathrm{X};\mathbf{C})\) , with support contained in U, such that \(\| h \| \leqslant 1\) and such that \(|\alpha_{k}(1 - \nu (h))| \leqslant \delta / 2\) for \(1 \leqslant k \leqslant n\) . The definition of U moreover shows that \(|(\alpha_{k} - f_{k}(x))h(x)| \leqslant \delta / 2\) for all \(x \in \mathrm{U}\) ; since \(\| \nu \| = 1\) and \(\operatorname{Supp}(\nu) \subset \mathrm{U}\) we therefore have \(|\nu ((\alpha_{k} - f_{k})h)| \leqslant \delta / 2\) and so, setting \(g = g_0h\) ,

\[
\left| f _ {k} \left(x _ {0}\right) - \mu \left(g f _ {k}\right) \right| \leqslant \delta \quad \text { for } 1 \leqslant k \leqslant n.
\]

This proves the proposition, since \(\| g\cdot \mu \| = \| (g_0h)\cdot \mu \| \leqslant \| g_0\cdot \mu \| = 1\)

COROLLARY. --- Let \(\mu\) be a measure on X. For a measure \(\nu\) on X to be in the vague closure of the set of measures \(g \cdot \mu\) , where \(g\) runs over \(\mathcal{K}(X; \mathbf{C})\) , it is necessary and sufficient that \(\operatorname{Supp}(\nu) \subset \operatorname{Supp}(\mu)\) .

For, \(\operatorname{Supp}(g\cdot\mu)\subset\operatorname{Supp}(\mu)\) by No.~3, Prop. 10; therefore the support of every vague limit of measures of the form \(g\cdot\mu\) is also contained in \(\operatorname{Supp}(\mu)\) (No.~2, Prop. 6). Conversely, if \(\operatorname{Supp}(\nu)\subset\operatorname{Supp}(\mu)\) then \(\nu\) is the vague limit of measures with finite support contained in \(\operatorname{Supp}(\mu)\) (Th. 1), hence is in the vague closure of the set of measures \(g\cdot\mu\) by Prop. 13.

